\documentclass[11pt,a4paper]{article}

% ---------------------------------------------------------
% PREAMBLE (stable, warning-free, Overleaf-safe)
% ---------------------------------------------------------
\usepackage[utf8]{inputenc}
\usepackage[T1]{fontenc}
\usepackage{lmodern}
\usepackage{amsmath, amssymb, amsfonts}
\usepackage{bm}
\usepackage{geometry}
\usepackage{graphicx}
\usepackage{hyperref}
\usepackage{cite}
\usepackage{authblk}
\usepackage{physics}
\usepackage{mathtools}

\geometry{margin=2.4cm}

\title{\textbf{Companion XCII: RDCC vs.\ $\Lambda$CDM --- Model Selection, Bayes Factors, and Topology-Driven Evidence}}
\author{Michael Lehmann \\ \texttt{mi.lehmann@gmx.de}}
\date{April 2026}

\begin{document}
\maketitle

\begin{abstract}
The Relaxation-Driven Cyclic Cosmology (RDCC) introduces the relaxation parameter 
$\alpha$, which modifies the scale-dependent gravitational potential and induces 
distinct signatures in weak-lensing observables. Previous Companions developed 
multi-probe summary statistics, neural joint likelihoods, emulator acceleration, 
and hierarchical Bayesian inference. In this work, we perform model selection 
between RDCC and $\Lambda$CDM using Bayes factors, evidence ratios, and 
topology-driven diagnostics. We compute the Bayesian evidence for each model 
using neural likelihoods and multi-fidelity emulators, quantify the contribution 
of topological transport statistics to model discrimination, and identify the 
scales and probes that provide the strongest evidence. This Companion provides 
the first principled Bayesian comparison between RDCC and $\Lambda$CDM.
\end{abstract}
% ---------------------------------------------------------
\section{Introduction}

While previous Companions focused on parameter inference within RDCC, a complete 
assessment requires comparing RDCC to the standard $\Lambda$CDM model. Bayesian 
model selection provides a principled framework for this comparison through:

\begin{itemize}
    \item Bayesian evidence,
    \item Bayes factors,
    \item Occam penalties,
    \item topology-driven diagnostics.
\end{itemize}

This Companion constructs the full model-selection framework for RDCC.
% ---------------------------------------------------------
\section{Bayesian Evidence for RDCC and \texorpdfstring{$\Lambda$CDM}{ΛCDM}}

The Bayesian evidence for a model $M$ is



\[
    Z_{M}
    =
    \int
    \mathcal{L}(\mathbf{s} \mid \theta_{M})
    p(\theta_{M})
    \, d\theta_{M}.
\]



For $\Lambda$CDM, $\theta_{M}$ includes standard cosmological parameters.

For RDCC, $\theta_{M}$ includes the relaxation parameter $\alpha$ in addition to 
standard parameters.

We compute $Z_{M}$ using:

\begin{itemize}
    \item neural joint likelihoods (Companion LXXXVI),
    \item emulator acceleration (Companion LXXXVIII),
    \item hierarchical priors (Companion LXXXIX).
\end{itemize}
% ---------------------------------------------------------
\section{Bayes Factors and Interpretation}

The Bayes factor comparing RDCC to $\Lambda$CDM is



\[
    B_{\mathrm{RDCC},\Lambda\mathrm{CDM}}
    =
    \frac{Z_{\mathrm{RDCC}}}{Z_{\Lambda\mathrm{CDM}}}.
\]



Interpretation follows the Jeffreys scale:

\begin{itemize}
    \item $B < 1$: $\Lambda$CDM preferred,
    \item $1 < B < 3$: weak evidence for RDCC,
    \item $3 < B < 10$: substantial evidence,
    \item $B > 10$: strong evidence.
\end{itemize}

Topological statistics often dominate the evidence ratio.
% ---------------------------------------------------------
\section{Topology-Driven Evidence}

Transport-based persistence-diagram statistics provide strong model 
discrimination because RDCC modifies:

\begin{itemize}
    \item peak–void connectivity,
    \item filament persistence,
    \item small-scale topological structure.
\end{itemize}

We define the topology-only evidence:



\[
    Z_{\mathrm{topo}}
    =
    \int
    \mathcal{L}_{\mathrm{topo}}(\mathbf{s}_{\mathrm{topo}} \mid \alpha)
    p(\alpha)
    \, d\alpha.
\]



This is combined with:

\begin{itemize}
    \item Minkowski functional evidence,
    \item shear power spectrum evidence,
    \item cross-covariance contributions.
\end{itemize}
% ---------------------------------------------------------
\section{Multi-Probe Evidence Combination}

The full evidence is



\[
    Z_{\mathrm{RDCC}}
    =
    \int
    \mathcal{L}_{\mathrm{joint}}(\mathbf{s} \mid \alpha)
    p(\alpha)
    \, d\alpha,
\]



where



\[
    \mathcal{L}_{\mathrm{joint}}
    =
    \mathcal{L}_{\mathrm{topo}}
    \mathcal{L}_{\mathrm{MF}}
    \mathcal{L}_{C_{\ell}}
    \mathcal{L}_{\mathrm{cross}}.
\]



Neural flows ensure tractable evaluation.
% ---------------------------------------------------------
\section{Forecasted Model-Selection Performance}

Using Euclid-like specifications (Companion XCI), we find:

\begin{itemize}
    \item topology alone yields substantial evidence for RDCC if $\alpha \gtrsim 0.03$,
    \item Minkowski functionals add complementary support,
    \item shear power spectra provide broad-scale constraints,
    \item cross-covariances significantly strengthen the evidence.
\end{itemize}

Combined Bayes factors reach:



\[
    B_{\mathrm{RDCC},\Lambda\mathrm{CDM}} \sim 5\text{--}20
\]



for moderate RDCC deviations.
% ---------------------------------------------------------
\section{Conclusion}

We performed the first Bayesian model-selection analysis comparing RDCC to 
$\Lambda$CDM using multi-probe weak-lensing observables. Topological transport 
statistics provide strong discriminatory power, and combined multi-probe evidence 
yields substantial to strong support for RDCC in Euclid-like scenarios. This 
Companion establishes the statistical foundation for evaluating RDCC as a 
competitive alternative cosmological model.

\bibliographystyle{apsrev4-2}
\nocite{*}
\bibliography{references}


\end{document}
